\section{Discussion}\label{sec:discussion}

\subsection{What the results are worth}
The five ``solved'' problems fall into three types. Two (18.46 and 13.19) are decided by an explicit object: a
group of order 128 (respectively $2^9$) whose relevant property is checked by a finite computation, plus a proof
that nothing smaller works (18.46) or a general construction (13.19). Two (21.32 and 2.78) are proved by short but
non-trivial arguments---a Baer-sum manipulation of central extensions, and a product formula together with an
explicit family of $p$-groups. One (16.60) is a reduction of the question to a classical theorem. In all cases the
proofs are complete and elementary, and the computations are reproducible from the repository. The agent was not
able to determine whether any of these answers already exist in the literature; the problems were all listed as
unsolved in the 21st edition, but 2.78 is sixty years old and 16.60 twenty, and it would not be surprising if
parts were known.

The partial results are of a different kind. For 16.14 and 20.37 the agent proved
theorems (a structure theorem for minimal counterexamples; a coset-union theorem obtained from the classical
perfect-matching property of vertex-transitive graphs) that reduce the problems to
finite checks in a range, and then carried out the checks. For 19.56 the reduction to Frattini extensions of
minimal simple groups is a clean statement, but the remaining step could not be done. For 20.21 the agent proved
two impossibility lemmas and searched well beyond the range of a direct search (its own bound: an example, if it exists, has order
at least 6144). In each case the report states precisely what is proved and what is
computed.

\subsection{Reliability}
Three mechanisms supported reliability. (1)~\emph{Independent verification of computational claims}: every
counterexample or construction was recomputed by a second method, and every 20.37 factorisation was rechecked by
brute-force multiplication in GAP. This caught a real bug (the coset-side error) and gave confidence in the others.
(2)~\emph{Self-review}: on the second day the agent re-read every report line by line, checked each step of the
proofs (recorded in the transcript), and corrected several imprecise statements (an unjustified ``more
generally'' clause in the 21.26 note, a vague ``chain'' claim in the 20.21 note, an overstated list of completely
settled groups in 20.37). (3)~\emph{Honest bookkeeping}: the verification table records ranges actually completed,
and the notes list literature dependencies that could not be checked.

The obvious weakness is the literature. The agent could search the web but could not read paywalled papers, and
did not attempt to find prior work on most of the problems it attacked. A human reviewer should therefore first
check whether the ``solved'' problems are already solved; the mathematics itself should be checkable in a few
hours per problem from the appendices.

A second weakness is that all mathematical claims were checked only by the same system that produced them. The
proofs in Appendix~\ref{app:proofs} are written so that a reader can verify each step; nothing depends on the
agent's authority.

\subsection{What worked and what did not}
The combination of a computer algebra system with an agent that can write and run code quickly is well suited to
this kind of problem list: many notebook problems ask whether a phenomenon occurs, and the SmallGroups library
answers such questions up to order 2000 in minutes. The agent's most productive pattern was ``compute first, then
prove'': the sharpness of Klyachko's bound, the class-2 construction for Gorchakov's problem, and the exclusion
strategy for Berkovich's problem all came from looking at data. Delegating the triage to sub-agents was
efficient (about 2300 problems read in an hour) but lossy (one report exceeded the output limit and part of it was
lost).

What did not work: attempts at general proofs where the problem is genuinely hard (the general case of 16.14, the Frattini lifting in 19.56, the case $I\cdot I'=K$ of 20.21); the agent recognised this fairly quickly each time and moved to
partial results. Two sweeps were started on a vacuous or wrong formulation and had to be redone. Long searches
sometimes ran without producing progress output, so their coverage could not be reported precisely, and one
(20.21, rank $\ge 5$ at order 512) turned out to be infeasible only after seven hours.

\subsection{Recommendations}
For anyone repeating such an experiment:
\begin{itemize}
\item Provide literature access (a searchable database of MathSciNet/zbMATH reviews or the notebook's archive of
solved problems); the largest uncertainty about the value of the results is their novelty.
\item Insist, as here, on independent verification of every computational claim; it is cheap and it works.
\item Make the agent write reports as it goes; the reports were the basis for the final self-review and for this
paper, and the git history is what allowed the timestamp error to be corrected.
\item Budget the final hours explicitly (the agent planned a consolidation phase but mis-scheduled it by misreading a
clock); a harness-provided deadline notification would have avoided this.
\item Prefer many small, well-instrumented sweeps with progress output over a few large ones.
\end{itemize}

\subsection{Outlook}
Several of the partial results suggest natural continuations: for 20.37, the cases $\PSL(3,4)$ and $\Sz(8)$ with
factor sizes 70 and 140 are small enough for an exhaustive (rather than randomised) search and would be the first
genuine test of the conjecture beyond subgroup-based constructions; for 16.14, the case $r=2$ could be closed with
the classification of 2-groups with three involutions, and the structure theorem may admit a purely theoretical
continuation through the 2-covering group; for 19.56, a uniform treatment of Frattini extensions of minimal simple
groups seems within reach of a specialist; for 2.78, the remaining question is whether every $k\ge 9$ is the
number of $\mathrm{IE}\bar n$-systems of some simple group, which depends on the arithmetic of the subgroup orders
of $\PSL(2,q)$.
